\documentclass[../../main.tex]{subfiles}
\graphicspath{{../../image/}{../../figure/}}
\begin{document}

\section{L$^{\infty}$ 解的精细结构}\label{section:ch06-Linfty解精细结构}

由\cref{theorem:ch06-6-2-6}，当初值 $u_{0}$ 在 $\mathbb{R}^{m}$ 上局部有界变差时，标量守恒律\eqref{align:ch06-6-1-1}的可容许解 $u$ 在上半空间上局部有界变差，从而继承\cref{section:ch01-BV函数理论概述}与\cref{section:ch01-BV解}中所述 $BV$ 函数的精细结构。

特别地，$u$ 的近似跳跃间断点组装在至多可数个余维一 $C^{1}$ 流形的并上。此外，$u$ 在任一余维一有向流形两侧具有（一般不同的）迹。然而，当 $u_{0}$ 仅属于 L$^{\infty}$ 时，上述结构一般丢失；考虑\eqref{align:ch06-6-1-1}为线性的情形即可见。另一方面，在\cref{section:ch06-动力论表述}（\cref{theorem:ch06-6-7-2}）中我们看到，流通量函数中的非线性可能对 L$^{\infty}$ 解施加磨光影响。

作为此现象的另一表现，此处将看到，当守恒律在某种意义下线性非退化时，仅属于 L$^{\infty}$ 的可容许解仍被赋予与 $BV$ 函数结构极为相似的精细结构。

就当前目的而言，空间变量与时间变量的区别无关紧要，故不妨回到\cref{chapter:ch01-平衡律}的表述与记号：把 $m$ 维空间与一维时间融合为 $k$ 维时空，$k=m+1$，用向量 $X$ 表示 $(x,t)$，其中 $X_{\alpha}=x_{\alpha}$（$\alpha=1,\cdots,m$），$X_{k}=t$。下文 $\operatorname{div}$ 表示 $\mathbb{R}^{k}$ 中作用于 $k$ 维行向量的散度算子。


在 $\mathbb{R}^{k}$ 的某开子集 $\mathscr{X}$ 上，考虑如下形式的标量平衡律
\begin{align}
\operatorname{div}G(u(X))=\nu_{G},
\label{align:ch06-6-8-1}
\end{align}
其中 $\nu_{G}$ 是 Radon 测度。若对 $G$ 的任一相伴 $Q$，
\begin{align}
\operatorname{div}Q(u(X))=\nu_{Q},
\label{align:ch06-6-8-2}
\end{align}
其中 $\nu_{Q}$ 是 $\mathscr{X}$ 上的 Radon 测度，则 $u\in L^{\infty}(\mathscr{X})$ 称为\eqref{align:ch06-6-8-1}的\emph{可容许解}。

回顾\cref{section:ch01-伴随平衡律}，相伴 $Q$ 与 $G$ 的关系为

\begin{align}
Q^{\prime}(u)=\eta^{\prime}(u)G^{\prime}(u),
\label{align:ch06-6-8-3}
\end{align}
其中 $\eta$ 是某标量值函数。

在\cref{section:ch06-可容许弱解稳定性}的框架下，$G_{k}(u)=u$、$Q_{k}(u)=\eta(u)$，且 $\nu_{G}=0$。

正如\cref{section:ch06-可容许弱解稳定性}所注，按\cref{definition:ch06-6-2-1}意义的任一可容许解 $u$，对任一相伴 $Q$（特别地，当 $Q_{k}$ 凸时为非正测度）都使分布 $\operatorname{div}Q(u)$ 成为测度，因而也是上述意义下的可容许解。


为排除线性组，引入下述概念（与\eqref{align:ch06-6-7-10}比较）：

\begin{definition}\label{definition:ch06-6-8-1}
若对每个 $N\in\mathbb{S}^{k-1}$，
\begin{align}
G^{\prime}(u)N\neq 0,\qquad\text{对几乎所有 }u\in(-\infty,\infty),
\label{align:ch06-6-8-4}
\end{align}
则称平衡律\eqref{align:ch06-6-8-1}是\emph{线性非退化}的。
\end{definition}

线性非退化标量平衡律可容许解的精细结构由下述定理描述。

\begin{theorem}\label{theorem:ch06-6-8-2}
设\eqref{align:ch06-6-8-1}线性非退化，$u$ 是 $\mathscr{X}$ 上的可容许解。则 $\mathscr{X}$ 是三个两两不相交子集 $\mathscr{C}$、$\mathscr{J}$、$\mathscr{I}$ 之并，具有如下性质：
\begin{enumerate}
\item[(a)] $\mathscr{C}$ 是 $u$ 的零平均振幅点集，即对 $\bar{X}\in\mathscr{C}$，
\begin{align}
\lim_{r\downarrow 0}\frac{1}{r^{k}}\int_{\mathscr{B}_{r}(\bar{X})}|u(X)-\bar{u}_{r}(\bar{X})|\,\mathrm{d}X=0,
\label{align:ch06-6-8-5}
\end{align}
其中 $\bar{u}_{r}(\bar{X})$ 表示 $u$ 在球 $\mathscr{B}_{r}(\bar{X})$ 上的平均。
\item[(b)] $\mathscr{J}$ 是可求长的，即它本质上被可数个嵌入 $\mathbb{R}^{k}$ 的 $C^{1}$ $(k-1)$ 维流形 $\{\mathscr{F}_{i}\}$ 之并所覆盖：$\mathscr{H}^{k-1}(\mathscr{J}\setminus\bigcup_{i}\mathscr{F}_{i})=0$。当 $\bar{X}\in\mathscr{J}\cap\mathscr{F}_{i}$ 时，$\mathscr{F}_{i}$ 在 $\bar{X}$ 处的法向即解释为 $\mathscr{J}$ 在 $\bar{X}$ 处的法向。函数 $u$ 在任一点 $\bar{X}\in\mathscr{J}$ 处具有不同的内向迹 $u_{-}$ 与外向迹 $u_{+}$（按\cref{definition:ch01-迹}意义）。

\item[(c)] $\mathscr{I}$ 的 $(k-1)$ 维 Hausdorff 测度为零：$\mathscr{H}^{k-1}(\mathscr{I})=0$。
\end{enumerate}
\end{theorem}

比较\cref{theorem:ch01-BV链式法则}与\cref{theorem:ch06-6-8-2}，可见可容许 L$^{\infty}$ 解与 $BV$ 函数在精细结构上惊人地相似。

然而读者应注意二者也有区别：在可容许 L$^{\infty}$ 解中，$\mathscr{C}$ 中点仅具有零平均振幅，而在 $BV$ 情形中它们是 Lebesgue 点。此外，若 $u$ 是\eqref{align:ch06-6-8-1}以 $\nu_{G}=0$ 的 $BV$ 解，则由\cref{theorem:ch01-BV伴随不等式}，对任一相伴 $Q$，$\nu_{Q}$ 集中于跳跃间断点集 $\mathscr{J}$ 上。

然而目前尚不知此重要性质能否推广到 L$^{\infty}$ 可容许解，除 $n=1$ 外；见\cref{section:ch11-注记}。


\cref{theorem:ch06-6-8-2}的证明冗长而技术性强，读者应参阅\cref{section:ch06-注记}所引文献。

尽管如此，此处仍简述若干关键要素。

$\mathscr{X}$ 上\eqref{align:ch06-6-8-1}的可容许 L$^{\infty}$ 解 $u$ 可由\cref{section:ch06-动力论表述}所讨论的动力论表述刻画。

在当前框架下，\eqref{align:ch06-6-7-2}取形式
\begin{align}
\sum_{\alpha=1}^{k}G_{\alpha}^{\prime}(v)\partial_{\alpha}\chi(v;u)=\partial_{v}\nu,
\label{align:ch06-6-8-6}
\end{align}
其中 $\chi$ 是\eqref{align:ch06-6-7-3}所定义的函数，$\nu$ 是 $\mathbb{R}\times\mathscr{X}$ 上的有界测度。注意与\cref{section:ch06-动力论表述}不同，此处测度 $\nu$ 不必非负，因为本节所采纳的可容许解概念更宽泛。


类比\eqref{align:ch06-6-7-6}，由某 $\eta$ 经\eqref{align:ch06-6-8-3}诱导的任一相伴 $Q$ 所对应的测度 $\nu_{Q}$ 与测度 $\nu$ 的关系为
\begin{align}
\nu_{Q}=-\int_{-\infty}^{\infty}\eta^{\prime\prime}(v)\,\mathrm{d}\nu(v;\cdot).
\label{align:ch06-6-8-7}
\end{align}

测度 $\nu$ 还通过下式确定\cref{theorem:ch06-6-8-2}中的“跳跃集” $\mathscr{J}$：
\begin{align}
\mathscr{J}=\Bigl{\{}X\in\mathscr{X}:\limsup_{r\downarrow 0}\frac{|\nu|(\mathbb{R}\times\mathscr{B}_{r}(X))}{r^{k-1}}>0\Bigr{\}},
\label{align:ch06-6-8-8}
\end{align}
其中 $|\nu|$ 表示 $\nu$ 的全变差测度。

$u$ 的精细结构的澄清是通过对任一点 $X\in\mathscr{X}$ 的邻域“爆破”实现的，即在 $X$ 附近以保持\eqref{align:ch06-6-8-6}不变的方式重标度 $u$ 与 $\nu$。线性非退化条件\eqref{align:ch06-6-8-4}结合输运方程\eqref{align:ch06-6-8-6}的速度平均估计，诱导出必要的紧性，使 $u$ 与 $\nu$ 在重标度下的极限 $u_{\infty}$、$\nu_{\infty}$ 存在且满足\eqref{align:ch06-6-8-6}。当 $X\notin\mathscr{J}$ 时，测度 $\nu_{\infty}$ 为零；反之，当 $X\in\mathscr{J}$ 时，$\nu_{\infty}$ 是 $\mathbb{R}$ 上测度与 $\mathscr{X}$ 上测度的张量积。正是通过研究具有此种特殊张量积结构的 $\nu$ 的\eqref{align:ch06-6-8-6}解，才确立了\cref{theorem:ch06-6-8-2}的论断。

用同样技巧可验证，线性非退化标量平衡律的可容许解与 $BV$ 函数共享另一重要性质：它们在余维一流形上具有单侧迹。

\begin{theorem}\label{theorem:ch06-6-8-3}
设 $u$ 是\eqref{align:ch06-6-8-1}在 $\mathbb{R}^{k}$ 中带边界 $\mathscr{B}$ 的 Lipschitz 子集 $\mathscr{X}$ 上的线性非退化平衡律的可容许解。假设对任一相伴 $Q$，\eqref{align:ch06-6-8-2}中的测度 $\nu_{Q}$ 在 $\mathscr{X}$ 上有限。则 $u$ 在 $\mathscr{B}$ 上具有强迹 $u_{\mathscr{B}}\in L^{\infty}(\mathscr{B})$。
\end{theorem}

强迹在 $L^{1}_{\mathrm{loc}}$ 中实现，大致如下：设 $\mathscr{B}$ 包含 $(k-1)$ 维超平面的紧子集 $\mathscr{P}$，其外单位法向为 $N$。则 $u_{\mathscr{B}}$ 在 $\mathscr{P}$ 上的限制由
\begin{align}
\operatorname*{ess\,lim}_{\tau\downarrow 0}\int_{\mathscr{P}}|u(X-\tau N)-u_{\mathscr{B}}(X)|\,\mathrm{d}\mathscr{H}^{k-1}(X)=0
\label{align:ch06-6-8-9}
\end{align}
刻画。

一般情形中，利用 $\mathbb{R}^{k}$ 上的 Lipschitz 变换把 $\mathscr{B}$ 的“片”映到超平面的“片” $\mathscr{P}$ 上，再使用上述刻画。

\cref{theorem:ch06-6-8-3}在标量守恒律边值问题理论中起重要作用，将在\cref{section:ch06-初边值问题}中看到。

\cref{theorem:ch06-6-8-3}的另一个重要推论如下。

\begin{corollary}\label{corollary:ch06-6-8-4}
设标量守恒律\eqref{align:ch06-6-1-1}线性非退化，$u$ 是上半空间上 Cauchy 问题\eqref{align:ch06-6-1-1}、\eqref{align:ch06-6-1-2}的 L$^{\infty}$ 弱解，在分布意义下对每个凸熵 $\eta$ 满足不等式\eqref{align:ch06-6-2-2}。则映射 $t\mapsto u(\cdot,t)$ 在 $L^{1}_{\mathrm{loc}}(\mathbb{R}^{m})$ 中对任一 $t\in[0,\infty)$ 强连续。
\end{corollary}

特别地，对线性非退化标量守恒律，Cauchy 问题的可容许解可仅由不等式组\eqref{align:ch06-6-2-2}刻画，而不必用更强的条件\eqref{align:ch06-6-2-3}。因而回到\cref{section:ch04-熵可容性条件}关于熵可容许性的讨论，可得出结论：对标量、线性非退化守恒律，集合 $\mathscr{F}$ 为空。


\end{document}
